% chapters/ch01-points-coords.tex

\chapter{点与坐标：第一次认知手术}

\begin{intuition}
	初学者看到极坐标 $(r,\theta)$ 时，大脑会自动把它们理解成“半径”和“角度”。
	这种直觉有用，但也可能遮蔽坐标本身的结构。
	本章先把“点”和“坐标表示”严格区分开来。
\end{intuition}

\section{流形上的点}

设 $M$ 是一个微分流形，$p \in M$ 是流形上的一个几何点。
在引入任何坐标之前，点 $p$ 本身就是客观存在的几何对象。

\begin{definition}[几何点]
	流形 $M$ 中的元素 $p$ 称为一个\emph{点}。它不依赖于任何坐标系。
\end{definition}

\section{坐标图与三层表示}

为了对点进行计算，我们需要将其“翻译”成数字。
标准微分几何中，坐标图通常直接写作：
\[
\varphi: U \subset M \to \mathbb{R}^n.
\]

本书为了教学清晰，引入一个辅助分解（\textbf{注意：这是教学框架，不是标准定义的强制改写}）：
\[
M \xrightarrow{\varphi} X \xrightarrow{\rho} \mathbb{R}^n
\]
其中：
\begin{itemize}
	\item $p \in M$ 是原始几何点；
	\item $p' = \varphi(p) \in X$ 是计算空间中的对应点；
	\item $\rho(p') = (x^1, \ldots, x^n)$ 是其坐标分量。
\end{itemize}

\begin{remark}
	当你阅读其他标准教材时，通常只看到 $\varphi: U \to \mathbb{R}^n$。
	本书的 $X$ 层是为了帮你建立“点 $\to$ 计算对象 $\to$ 数字”的认知层次，
	不要误以为所有微分几何教材都必须包含这个中间层。
\end{remark}

\section{极坐标示例}

考虑 $\mathbb{R}^2$ 中的点 $p$。
极坐标映射 $\varphi$ 将其对应到 $(r,\theta)$，其中：
\[
x = r\cos\theta,\quad y = r\sin\theta.
\]

\begin{example}
	点 $p$ 在极坐标下表示为 $(r,\theta) = (1, \pi/4)$。
	这里 $(1, \pi/4)$ 是坐标表示，不是点本身。
\end{example}

\begin{exercise}
	在球坐标系中，点 $p$ 的坐标表示为 $(r,\theta,\phi)$。
	写出从球坐标到笛卡尔坐标的映射 $\rho$ 的具体公式。
\end{exercise}